\section{Nghiên cứu tiên tiến liên quan}
\label{sec:related-research}

\subsection{Bài báo được khảo sát}

Phần này khảo sát bài báo của Wang, Huang, Rudin và Shaposhnik,
\textit{Understanding How Dimension Reduction Tools Work: An Empirical
Approach to Deciphering t-SNE, UMAP, TriMap, and PaCMAP for Data
Visualization}, công bố trên \textit{Journal of Machine Learning Research}
\cite{wang2021pacmap}. Bài báo nghiên cứu các thuật toán giảm chiều phục
vụ trực quan hóa, đặc biệt trong tình huống dữ liệu chiều cao được nhúng
xuống hai chiều. Ngoài phân tích t-SNE, UMAP và TriMap, tác giả đề xuất
thuật toán PaCMAP (\textit{Pairwise Controlled Manifold Approximation
Projection}).

Nội dung này nối trực tiếp với các phương pháp đã trình bày ở các chương
trước. PCA cung cấp phép chiếu tuyến tính và thường là khởi tạo cho các
thuật toán hiện đại; Isomap, Laplacian eigenmaps và LLE biểu diễn quan điểm
dữ liệu nằm trên manifold thông qua quan hệ lân cận; KPCA cho thấy cách tạo
embedding phi tuyến từ độ tương tự. Bài báo tiếp tục câu hỏi đó trong bối
cảnh trực quan hóa: một embedding đẹp cần giữ thông tin lân cận nào và giữ
bao nhiêu cấu trúc toàn cục.

\subsection{Bài toán nghiên cứu}

Một scatter plot hai chiều có thể khiến người đọc diễn giải khoảng cách
giữa cụm, kích thước cụm hoặc các khoảng trống như cấu trúc thật của dữ
liệu gốc. Tuy nhiên, khi giảm từ số chiều lớn xuống hai chiều, không thể
bảo toàn đồng thời mọi khoảng cách. t-SNE và UMAP thường làm nổi bật láng
giềng cục bộ, trong khi cấu trúc giữa các nhóm xa nhau có thể bị biến dạng.
Ngược lại, một phương pháp quá ưu tiên cấu trúc toàn cục có thể làm mờ các
cụm nhỏ có ý nghĩa.

Từ vấn đề này, bài báo đặt ra ba câu hỏi chính:
\begin{enumerate}
  \item Vì sao các thuật toán phổ biến tạo ra những embedding khác nhau,
  dù cùng được dùng cho trực quan hóa?
  \item Thành phần nào của hàm mất mát, đồ thị hoặc phép khởi tạo quyết định
  khả năng giữ cấu trúc cục bộ và toàn cục?
  \item Có thể thiết kế một thuật toán cân bằng hai mục tiêu đó tốt hơn hay
  không?
\end{enumerate}

\subsection{Đóng góp chính của bài báo}

\paragraph{Phân tích hàm mất mát bằng lực hút và lực đẩy.}
Các tác giả diễn giải thuật toán trực quan hóa thông qua lực hút giữa những
điểm cần ở gần nhau và lực đẩy giữa những điểm không nên chồng lên nhau.
Thay vì chỉ so sánh công thức, bài báo đưa ra biểu diễn ``rainbow figure''
để quan sát độ lớn gradient theo khoảng cách trong không gian gốc và không
gian nhúng. Từ phân tích thực nghiệm này, tác giả tổng hợp sáu nguyên tắc
cho một hàm mất mát có khả năng giữ cấu trúc cục bộ tốt; chẳng hạn điểm
láng giềng bị đặt xa cần chịu lực hút đủ mạnh, còn điểm không lân cận đã ở
xa không nên tiếp tục chi phối quá mức quá trình tối ưu
\cite{wang2021pacmap}.

\paragraph{Vai trò của đồ thị và phép khởi tạo.}
Bài báo chỉ ra rằng chỉ thu hút các láng giềng gần nhất và đẩy các điểm xa
khi chúng va vào nhau chưa đủ để phục hồi cấu trúc toàn cục. Các cặp điểm
ở khoảng cách trung gian cung cấp thông tin quan trọng về bố cục tổng thể.
Ngoài ra, hiệu quả giữ cấu trúc toàn cục của TriMap phụ thuộc đáng kể vào
khởi tạo PCA; điều này nhấn mạnh rằng kết quả embedding không chỉ do hàm
mất mát, mà còn do trạng thái ban đầu và thang đo của khởi tạo.

\paragraph{Thuật toán PaCMAP.}
PaCMAP xây dựng ba loại cặp điểm:
\begin{itemize}
  \item \textit{neighbor pairs}: các cặp láng giềng dùng để duy trì cấu
  trúc cục bộ;
  \item \textit{mid-near pairs}: các cặp ở khoảng cách trung gian giúp định
  hình cấu trúc toàn cục;
  \item \textit{further pairs}: các cặp xa giúp tránh việc embedding co cụm
  không hợp lý.
\end{itemize}
Thuật toán tối ưu theo ba giai đoạn. Giai đoạn đầu tăng trọng số của
\textit{mid-near pairs} để dựng bố cục lớn; giai đoạn giữa cân bằng các
loại cặp; giai đoạn cuối giảm vai trò của cặp trung gian và tinh chỉnh
quan hệ lân cận. Theo các thí nghiệm trên dữ liệu ảnh, văn bản, dữ liệu
sinh học và dữ liệu tổng hợp trong bài báo, PaCMAP đạt cân bằng tốt giữa
đánh giá cục bộ và toàn cục, đồng thời có thời gian chạy cạnh tranh.

\begin{table}[H]
  \centering
  \caption{Góc nhìn của Wang et al. về các thuật toán trực quan hóa giảm chiều.}
  \label{tab:related-methods}
  \begin{tabular}{p{0.18\textwidth}p{0.32\textwidth}p{0.39\textwidth}}
    \toprule
    Phương pháp & Thông tin ưu tiên & Nhận xét liên quan trong bài báo \\
    \midrule
    t-SNE & Láng giềng cục bộ & Cho cụm cục bộ rõ, nhưng bố cục toàn cục
    không nên được diễn giải trực tiếp. \\
    UMAP & Đồ thị lân cận mờ & Có dạng lực phù hợp với các nguyên tắc giữ
    cục bộ được phân tích. \\
    TriMap & Quan hệ bộ ba & Khả năng giữ toàn cục chịu ảnh hưởng lớn từ
    khởi tạo PCA. \\
    PaCMAP & Cặp gần, trung gian và xa & Dùng lịch trọng số theo giai đoạn
    để cân bằng cục bộ và toàn cục. \\
    \bottomrule
  \end{tabular}
\end{table}

\subsection{Liên hệ với lý thuyết giảm chiều đã trình bày}

\begin{table}[H]
  \centering
  \caption{Quan hệ giữa bài báo khảo sát và nội dung lý thuyết của báo cáo.}
  \label{tab:related-theory}
  \begin{tabular}{p{0.20\textwidth}p{0.70\textwidth}}
    \toprule
    Nội dung chương & Liên hệ với Wang et al. \\
    \midrule
    PCA & PCA là baseline tuyến tính giữ phương sai/toàn cục và còn được
    dùng làm khởi tạo cho một số embedding phi tuyến; vì vậy chất lượng
    khởi tạo có thể ảnh hưởng đến hình trực quan cuối cùng. \\
    KPCA & KPCA và các thuật toán trong bài báo đều tạo embedding phi tuyến,
    nhưng KPCA tối ưu theo phổ của kernel, còn PaCMAP tối ưu lực trên các
    cặp điểm được lựa chọn; hai cách tiếp cận không bảo toàn cùng một loại
    hình học. \\
    Isomap & Isomap giữ khoảng cách geodesic trên đồ thị để mô hình hóa hình
    học toàn cục của manifold. PaCMAP cũng dùng quan hệ láng giềng, nhưng
    bổ sung cặp trung gian và xa để phục vụ trực quan hóa cân bằng hơn. \\
    Laplacian/LLE & Các phương pháp này nhấn mạnh cấu trúc địa phương trên
    đồ thị; phân tích của bài báo giải thích vì sao chỉ thông tin địa phương
    có thể làm mất bố cục toàn cục. \\
    JL & Bổ đề JL đưa ra bảo đảm méo khoảng cách cho chiếu ngẫu nhiên ở số
    chiều đủ lớn; trái lại, bài toán trực quan hóa xuống đúng hai chiều
    chấp nhận biến dạng và phải chọn loại cấu trúc cần ưu tiên. \\
    \bottomrule
  \end{tabular}
\end{table}

Như vậy, bài báo không thay thế các kết quả phổ hoặc bảo đảm toán học của
chương, mà bổ sung một góc nhìn thiết kế thuật toán: khi mục đích là quan
sát embedding hai chiều, việc kiểm soát cặp điểm và gradient quyết định
cách đánh đổi giữa cấu trúc gần và bố cục xa.

\subsection{[MỞ RỘNG] Đối chiếu với cài đặt của nhóm}

Trong mã nguồn của đồ án, lớp \texttt{TSNE} cung cấp giao diện thống nhất
cho \texttt{sklearn.manifold.TSNE}, còn lớp \texttt{UMAP} được cài đặt với
các bước xây đồ thị $k$-láng giềng, tính trọng số fuzzy, khởi tạo phổ và
tối ưu với mẫu âm. Notebook thực nghiệm dùng mẫu cố định gồm 500 ảnh MNIST
và \texttt{random\_state = 42}; nhãn chỉ được dùng để đánh giá sau khi
nhúng. Hai chỉ số được báo cáo là \textit{trustworthiness@12}, đo mức bảo
toàn láng giềng, và độ chính xác KNN năm láng giềng trên embedding hai
chiều.

\begin{table}[H]
  \centering
  \caption{Thực nghiệm mở rộng trên 500 mẫu MNIST trong notebook.}
  \label{tab:bonus-tsne-umap}
  \begin{tabular}{lcc}
    \toprule
    Phương pháp & Trustworthiness@12 & KNN accuracy@5 \\
    \midrule
    t-SNE, perplexity $=30$ & \textbf{0.9291} & \textbf{0.800} \\
    t-SNE, perplexity $=10$ & 0.9284 & 0.798 \\
    t-SNE, perplexity $=50$ & 0.9282 & 0.782 \\
    t-SNE, perplexity $=5$  & 0.9195 & 0.766 \\
    Isomap & 0.7525 & 0.424 \\
    PCA & 0.7314 & 0.316 \\
    UMAP tự cài & 0.5006 & 0.112 \\
    \bottomrule
  \end{tabular}
\end{table}

Kết quả này phù hợp với cảnh báo của bài báo ở hai điểm. Thứ nhất, t-SNE
nhạy với tham số: chỉ thay \texttt{perplexity} đã làm hai chỉ số thay đổi.
Thứ hai, không nên suy diễn từ tên thuật toán rằng một embedding luôn tốt
hơn thuật toán khác. Cài đặt UMAP hiện tại cho kết quả thấp, vì vậy chỉ
được ghi nhận như thử nghiệm mở rộng cần hiệu chỉnh thêm; báo cáo không
dùng kết quả đó để khẳng định UMAP kém về bản chất hay t-SNE luôn vượt trội.

\subsection{Xu hướng, hạn chế và câu hỏi mở}

Từ bài báo có thể nhận xét xu hướng nghiên cứu giảm chiều phục vụ trực
quan hóa đang chuyển từ một tiêu chuẩn duy nhất, như phương sai của PCA
hoặc khoảng cách geodesic của Isomap, sang thiết kế mục tiêu lai: cấu trúc
cục bộ phải dễ quan sát nhưng bố cục toàn cục không được tùy tiện. Đồng
thời, cộng đồng quan tâm hơn đến việc giải thích hàm mất mát, đánh giá bằng
nhiều chỉ số và khảo sát độ nhạy tham số thay vì chỉ đưa ra một hình ảnh
embedding.

Bài báo chủ yếu đưa ra lập luận thiết kế và bằng chứng thực nghiệm cho
trực quan hai chiều; chất lượng hình thu được vẫn phụ thuộc tập dữ liệu,
siêu tham số, khởi tạo và chỉ số đánh giá. Một hướng kiểm chứng tiếp theo
cho đồ án là cài đặt PaCMAP hoặc dùng hiện thực chuẩn, rồi so sánh trên
cùng dữ liệu với PCA, Isomap, t-SNE và UMAP bằng cả thước đo cục bộ lẫn
thước đo toàn cục như độ chính xác triplet mà Wang et al. sử dụng.
